% =====================================================================
% Appendix section: Effect of the ionizing SED shape on the Nitrogen lines
% Paste this block into metal_gradient.tex *after* the existing
% "Effect of microturbulence on the Nitrogen lines" appendix and *before*
% the \bibliography line.  Copy the figure
%   plots/nitrogen/aox_line_ratios.png
% from the metallicity_gradient repository into figures/observation/ so
% the \includegraphics path below resolves.
% =====================================================================

\section{Impact of ionizing SED shape on the Nitrogen abundance}
\label{sec:appendix_aox}

A second candidate explanation for the strong $\rm N\,V/C\,IV$ relative to
$\rm N\,IV]/C\,IV$ and $\rm N\,III]/C\,IV$ at a single nitrogen abundance
is a harder ionizing SED, which selectively boosts $\rm N\,V$ (which
requires $h\nu > 77~\rm eV$ photons to ionize $\rm N^{3+}\rightarrow
N^{4+}$) without similarly boosting the lower-ionization nitrogen lines.
We test this hypothesis directly by varying the X-ray-to-UV slope
parameter $\alpha_{\rm ox}$ in the input ionizing continuum.

We use \textsc{Cloudy} version 23.01 \citep{Ferland2017} with the built-in
\texttt{AGN} continuum command, in which the incident SED is parameterized
as a UV ``big blue bump'' modified by an X-ray power law,
\begin{equation}
f_{\nu} = \nu^{\alpha_{\rm uv}}\,
         \exp\!\left(-h\nu/kT_{\rm BB}\right)\,
         \exp\!\left(-kT_{\rm IR}/h\nu\right)
       + a\,\nu^{\alpha_{x}}.
\label{eq:agn_sed}
\end{equation}
The constant $a$ is set internally by the requested $\alpha_{\rm ox}$,
defined in the standard way as the slope of the line connecting the
monochromatic luminosities at $2~\rm keV$ and $2500~\rm \AA$
\citep{Tananbaum1979}.  We fix $T_{\rm BB} = 10^{5}~\rm K$,
$\alpha_{\rm uv} = -0.5$, and $\alpha_{x} = -1.0$, and vary
$\alpha_{\rm ox}$ from $-2.0$ to $-1.0$ in steps of $0.1$.  This range
brackets the entire physical distribution of $\alpha_{\rm ox}$ observed in
luminous type-1 quasars
\citep{Strateva2005, Steffen2006, Just2007, LussoRisaliti2016}, whose
median lies near $\alpha_{\rm ox} \simeq -1.5$.

For each SED we run a single-zone Cloudy model with hydrogen density
$\log(n_{\rm H}/{\rm cm^{-3}}) = 10$, ionizing photon flux
$\log[\phi({\rm H})/(\rm cm^{-2}\,s^{-1})] = 19$, microturbulence
$v_{\rm turb} = 100~\rm km\,s^{-1}$, and a stopping column density
$\log(N_{\rm H}/{\rm cm^{-2}}) = 24$, varying the $\\alpha$-element abundance on
a logarithmic grid from $\log Z_{\alpha}/Z_{\odot} = -1$ to $+1$ in steps of
$0.2~\rm dex$.  These conditions are characteristic of the BLR clouds
that dominate the emission of the $\rm N\,V$, $\rm N\,IV]$, $\rm N\,III]$,
and $\rm C\,IV$ lines.

Figure~\ref{fig:aox_lineratios} shows the resulting line ratios versus
$\alpha_{\rm ox}$, color-coded by the $\\alpha$-element abundance.  The top row
displays the three nitrogen-to-$\rm C\,IV$ ratios that drive the
abundance discrepancy ($\rm N\,V$, $\rm N\,IV]$, $\rm N\,III]$); the
bottom row shows the low-ionization sanity-check ratios
$\rm Mg\,II/C\,IV$, $\rm C\,III]/C\,IV$, and $\rm C\,II]/C\,IV$.  Observed
broad-line composite values from \citet{VandenBerk2001} (red) and
\citet{constantin2003} (blue) are overlaid as horizontal bands.

To define what range of $\alpha_{\rm ox}$ is consistent with the
\emph{low-ionization} part of the BLR spectrum at solar $\\alpha$-element abundance, we
compute the $\alpha_{\rm ox}$ values at which the model curves at
$\log Z_{\alpha}/Z_{\odot} = 0$ simultaneously match the
\citet{VandenBerk2001} bands for $\rm Mg\,II/C\,IV$, $\rm C\,III]/C\,IV$, and
$\rm C\,II]/C\,IV$ within their $1\sigma$ measurement errors.  This
constraint pins $\alpha_{\rm ox}$ to a narrow range,
$\alpha_{\rm ox} \in [-1.67,\,-1.63]$ at $\log Z_{\alpha}/Z_{\odot} = 0$, shown
as the grey vertical strip in every panel.  Reassuringly, this
``low-ion-consistent'' $\alpha_{\rm ox}$ lies well within the
observationally established quasar distribution.

Within this strip the three nitrogen line ratios at solar $\\alpha$-element abundance
deviate from the observed values \emph{in opposite directions}:
$\rm N\,V/C\,IV$ is under-predicted by a factor of $\sim2$--$3$ relative
to \citet{VandenBerk2001} (and by a factor of $\sim7$ relative to
\citet{constantin2003}), whereas $\rm N\,IV]/C\,IV$ and
$\rm N\,III]/C\,IV$ are simultaneously \emph{over-predicted} by factors
of $\sim3$--$5$.  No combination of $\alpha_{\rm ox}$ within the
physically plausible quasar range can simultaneously bring all three
nitrogen ratios into agreement with the observed composites at a single
nitrogen abundance.  Even at $\alpha_{\rm ox} = -1.0$, which is harder
than essentially all observed type-1 quasars, the model
$\rm N\,V/C\,IV$ at $\log Z_{\alpha}/Z_{\odot} = 0$ remains a factor of
$\sim 2$--$3$ below the \citet{VandenBerk2001} value, while $\rm N\,III]/C\,IV$
falls below its observed band only at very subsolar $\\alpha$-element abundance.

We therefore conclude that the discrepancy between the nitrogen
abundance inferred from $\rm N\,V$ and from $\rm N\,IV]$/$\rm N\,III]$
cannot be removed by adjusting the shape of the ionizing SED at fixed
metallicity.  This residual discrepancy further supports our
interpretation that an enhancement of the nitrogen abundance over the
scaled-solar value is required to reproduce the observed $\rm N\,V$
strength in the BLR.  The remaining (sub-dex) systematic uncertainty in
the inferred $\alpha_{\rm ox}$, which translates into a few-tenths-dex
uncertainty in the inferred N abundance, is propagated into the error
budget of our nitrogen-gradient determination.

\begin{figure*}
\centering
\includegraphics[width=0.95\textwidth]{figures/observation/aox_line_ratios.png}
\caption{Cloudy single-zone model line ratios as a function of
$\alpha_{\rm ox}$, color-coded by $\\alpha$-element abundance
$\log Z_{\alpha}/Z_{\odot}$.  Top row: nitrogen-to-$\rm C\,IV$ ratios
($\rm N\,V$, $\rm N\,IV]$, $\rm N\,III]$); bottom row: low-ionization
checks ($\rm Mg\,II$, $\rm C\,III]$, $\rm C\,II]$).  Red and blue dashed
lines/bands show the broad-line composite values of \citet{VandenBerk2001} and
\citet{constantin2003} respectively.  The grey vertical strip marks the
$\alpha_{\rm ox}$ range that simultaneously reproduces the
low-ionization $\rm Mg\,II$, $\rm C\,III]$, and $\rm C\,II]$ ratios at
solar $\\alpha$-element abundance within the \citet{VandenBerk2001} $1\sigma$ bands.  Within
that strip the three nitrogen ratios deviate from the observed values
in opposite senses, demonstrating that no choice of SED hardness can
simultaneously fit all three at a single nitrogen abundance.}
\label{fig:aox_lineratios}
\end{figure*}
